\section{Results From Prior NSF Support}

\BfPara{Award \# 2335700: Education DCL: EAGER: Developing Experiential Cybersecurity and Privacy Training for AI Practitioners} (PI: Abuhamad, 11/01/2023 -- 10/31/2025, \$299,905.00). 
\textbf{Intellectual Merit:} The project establishes experiential cybersecurity training for AI professionals. The goal
is to raise AI workers' awareness of security and privacy risks by developing and evaluating a comprehensive 12-workshop experiential training program. The workshop series provides the knowledge and skills needed to build AI systems that are not only technically sound from an AI perspective, but also secure, ethical, and privacy-preserving. 
\textbf{Broader Impacts:}
The completion of the program will result in a workforce proficient in the use of AI with an understanding of its cybersecurity issues. The program's materials and outcomes will be shared through various channels to promote responsible AI practices. Educational resources, case studies, and research papers will be available on a dedicated website. %Findings will also be published in reputable conferences and journals.




\BfPara{Award \# 2336386: CyberCorps Scholarship for Service: CyberRamblers Bolster National Security} (Co-PI: Abuhamad, 01/01/2024 --  12/31/2028, \$3,803,006.00). 
\textbf{Intellectual Merit:}
The CyberRamblers program enhances the interdisciplinary activities of the Loyola Center for Cybersecurity. Scholars in this program engage in interdisciplinary research and education, which fosters critical thinking, independence, teamwork, and technical knowledge essential for a cybersecurity career. The B.S. in Cybersecurity curriculum emphasizes a hands-on approach to learning, enabling students to apply concepts and confidently utilize the latest tools.
\textbf{Broader Impacts:}
Twenty undergraduate students will be trained through a program aimed at preparing them for careers in cybersecurity, addressing the high demand for professionals in this field. The program's goal includes having 50\% of participants from underrepresented groups. This initiative will benefit society by enhancing national security and protecting various critical sectors such as energy, commerce, finance, education, healthcare, manufacturing, and agriculture against cyber-attacks.